\documentclass{article}

% if you need to pass options to natbib, use, e.g.:
%     \PassOptionsToPackage{numbers, compress}{natbib}
% before loading neurips_2026

% The authors should use one of these tracks.
% Before accepting by the NeurIPS conference, select one of the options below.
% 0. "default" for submission
\usepackage{neurips_2026}
% the "default" option is equal to the "main" option, which is used for the Main Track with double-blind reviewing.
% 1. "main" option is used for the Main Track
%  \usepackage[main]{neurips_2026}
% 2. "position" option is used for the Position Paper Track
%  \usepackage[position]{neurips_2026}
% 3. "eandd" option is used for the Evaluations & Datasets Track
 % \usepackage[eandd]{neurips_2026}
 % if you need to opt-in for a single-blind submission in the E&D track:
 %\usepackage[eandd, nonanonymous]{neurips_2026}
% 4. "creativeai" option is used for the Creative AI Track
%  \usepackage[creativeai]{neurips_2026}
% 5. "sglblindworkshop" option is used for the Workshop with single-blind reviewing
 % \usepackage[sglblindworkshop]{neurips_2026}
% 6. "dblblindworkshop" option is used for the Workshop with double-blind reviewing
%  \usepackage[dblblindworkshop]{neurips_2026}

% After being accepted, the authors should add "final" behind the track to compile a camera-ready version.
% 1. Main Track
 % \usepackage[main, final]{neurips_2026}
% 2. Position Paper Track
%  \usepackage[position, final]{neurips_2026}
% 3. Evaluations & Datasets Track
 % \usepackage[eandd, final]{neurips_2026}
% 4. Creative AI Track
%  \usepackage[creativeai, final]{neurips_2026}
% 5. Workshop with single-blind reviewing
%  \usepackage[sglblindworkshop, final]{neurips_2026}
% 6. Workshop with double-blind reviewing
%  \usepackage[dblblindworkshop, final]{neurips_2026}
% Note. For the workshop paper template, both \title{} and \workshoptitle{} are required, with the former indicating the paper title shown in the title and the latter indicating the workshop title displayed in the footnote.
% For workshops (5., 6.), the authors should add the name of the workshop, "\workshoptitle" command is used to set the workshop title.
% \workshoptitle{WORKSHOP TITLE}

% "preprint" option is used for arXiv or other preprint submissions
 % \usepackage[preprint]{neurips_2026}

% to avoid loading the natbib package, add option nonatbib:
%    \usepackage[nonatbib]{neurips_2026}

% Standard package includes
\usepackage{times}
\usepackage{latexsym}

% For proper rendering and hyphenation of words containing Latin characters (including in bib files)
\usepackage[T1]{fontenc}
% For Vietnamese characters
% \usepackage[T5]{fontenc}
% See https://www.latex-project.org/help/documentation/encguide.pdf for other character sets

% This assumes your files are encoded as UTF8
\usepackage[utf8]{inputenc}

% This is not strictly necessary, and may be commented out,
% but it will improve the layout of the manuscript,
% and will typically save some space.
\usepackage{microtype}

% This is also not strictly necessary, and may be commented out.
% However, it will improve the aesthetics of text in
% the typewriter font.
\usepackage{inconsolata}

%Including images in your LaTeX document requires adding
%additional package(s)
\usepackage{graphicx}

% If the title and author information does not fit in the area allocated, uncomment the following
%
%\setlength\titlebox{<dim>}
%
% and set <dim> to something 5cm or larger.

\usepackage{amsmath}
\usepackage{amssymb}
\usepackage{booktabs}
\usepackage{multirow}
\usepackage{xcolor}
\usepackage{subcaption}
\usepackage{float}
\newcommand{\method}{HyperSafe}
\newcommand{\ssn}{SSN}
\newcommand{\lam}{\lambda}
\newcommand{\remark}[1]{{\color{red}\textbf{[#1]}}}

% Note. For the workshop paper template, both \title{} and \workshoptitle{} are required, with the former indicating the paper title shown in the title and the latter indicating the workshop title displayed in the footnote. 
\title{Inference-Time Safety Recovery for Fine-Tuned Language Models \\via Activation-Conditioned Hypernetworks}


% The \author macro works with any number of authors. There are two commands
% used to separate the names and addresses of multiple authors: \And and \AND.
%
% Using \And between authors leaves it to LaTeX to determine where to break the
% lines. Using \AND forces a line break at that point. So, if LaTeX puts 3 of 4
% authors names on the first line, and the last on the second line, try using
% \AND instead of \And before the third author name.


\author{%
  David S.~Hippocampus\thanks{Use footnote for providing further information
    about author (webpage, alternative address)---\emph{not} for acknowledging
    funding agencies.} \\
  Department of Computer Science\\
  Cranberry-Lemon University\\
  Pittsburgh, PA 15213 \\
  \texttt{hippo@cs.cranberry-lemon.edu} \\
  % examples of more authors
  % \And
  % Coauthor \\
  % Affiliation \\
  % Address \\
  % \texttt{email} \\
  % \AND
  % Coauthor \\
  % Affiliation \\
  % Address \\
  % \texttt{email} \\
  % \And
  % Coauthor \\
  % Affiliation \\
  % Address \\
  % \texttt{email} \\
  % \And
  % Coauthor \\
  % Affiliation \\
  % Address \\
  % \texttt{email} \\
}


\begin{document}


\maketitle


% HyperSafe
\begin{abstract}
Safety alignment in large language models is fragile under fine-tuning, as even benign adaptation can increase harmful compliance.
Existing defenses rely on per-checkpoint interventions (either retraining or modifying the model weights) or on model-agnostic classifiers that do not capture model-specific shifts, forcing a trade-off between deployment cost and recipe-specific protection.
We propose HyperSafe, a framework that generates a model-specific safety classifier at inference time.
Our approach builds on the observation that fine-tuning induces characteristic shifts in intermediate representations, which we use as an activation fingerprint of the fine-tuned model.
A hypernetwork maps these fingerprints to the parameters of a \textit{Safe Side Network} (\ssn{}) that classifies prompts as safe or harmful.
At deployment, a single forward pass produces an \ssn{} from a small set of calibration prompts, without gradient updates or model modification.
We evaluate HyperSafe on two model families (Qwen2-7B and LLaMA-3-8B) across three safety benchmarks (BeaverTails, AdvBench, HEx-PHI).
HyperSafe reduces harmful response rates from 19--31\% to under 1\% on every held-out checkpoint, while keeping downstream task accuracy within one percentage point of the fine-tuned baseline on average.
\end{abstract}

\section{Introduction}
\label{sec:introduction}
%TODO
% 1. Check Text
% 2. Check references
% 3. Adjust contributions
% 4. Include figure


Large language models (LLMs) are commonly released with safety alignment that reduces harmful outputs, typically via instruction tuning and preference optimization. However, a growing body of evidence shows that this alignment can be brittle: downstream fine-tuning, including seemingly benign adaptation, may degrade refusal behavior and increase vulnerability to jailbreaks and harmful fine-tuning attacks \cite{qi2024finetuning,qi2025tokensdeep,fraser2025finetuninglowers,huang2024lisa}. This fragility poses a practical challenge: the open ecosystem increasingly relies on fine-tuning, adapters, and frequent checkpointing, yet safety guarantees do not reliably transfer across these model variants.

A natural response is to add guardrails or classifiers that detect harmful prompts. Prior work explores lightweight safety classifiers and adversarially robust prompt shields \cite{kim2024robust}, as well as activation-based detectors that monitor intermediate representations to flag attacks \cite{jiang2025hiddendetect}. Separately, post-fine-tuning alignment methods aim to restore safety by additional training or alignment steps \cite{huang2024lisa}. While effective in some regimes, these approaches often require repeated retraining, access to substantial safety data, or careful tuning per target model, limiting scalability.
%when many fine-tuned models must be monitored or deployed.

We take a different perspective: prompt safety is model-dependent. The same prompt can be handled safely by one fine-tuned model yet elicit unsafe compliance from another. This suggests learning a model-conditioned safety mechanism, i.e., estimating a safety score given both the prompt and the particular fine-tuned model. Our central goal is to scale safety assessment to unseen fine-tuned checkpoints without retraining a new safety classifier each time.

To this end, we propose activation-conditioned safety head generation. For each fine-tuned LLM, we extract intermediate activations while processing inputs and train a small side network (the ``safety head'') that maps these activations to a scalar safety score. We then train a hypernetwork \cite{chauhan2024brief} that takes the fine-tuned model's activations as input and outputs the weights of the safety head. At deployment, given a new fine-tuned model, we extract activations and use the hypernetwork to instantiate a model-specific safety head without additional safety training.

A key design choice is to generate the safety head in a layer-wise, factorized manner. Directly mapping all activations to all safety-head weights in one step yields an ill-conditioned, high-dimensional regression problem. Instead, we exploit the compositional structure of the safety head and predict each layer's parameters conditioned on corresponding LLM activations, reducing output dimensionality, improving identifiability, and stabilizing training. 
% Optionally, we incorporate behavioral probe fingerprints as auxiliary context to improve generalization across fine-tuning recipes and to disambiguate cases where activations alone are insufficient.

Overall, our method offers a scalable route to post-fine-tuning safety assessment: rather than re-aligning every new checkpoint, we amortize the cost by learning a generator that instantiates a lightweight, model-specific safety head from activations. This bridges lines of work on prompt safety detection \cite{kim2024robust,jiang2025hiddendetect}, alignment fragility under fine-tuning \cite{qi2024finetuning,qi2025tokensdeep,fraser2025finetuninglowers}, and hypernetwork-based weight generation \cite{chauhan2024brief}.

\paragraph{Contributions.}
\begin{itemize}
    % \item We introduce a model-conditioned safety formulation for fine-tuned LLMs, emphasizing that prompt safety depends on the particular fine-tuned checkpoint.
    \item We propose activation-conditioned hypernetwork safety head generation: a hypernetwork maps intermediate LLM activations to the weights of a lightweight side safety network, enabling zero-shot instantiation of safety heads for unseen fine-tuned models.
    \item We develop a layer-wise, factorized hypernetwork that predicts each safety-head layer from corresponding LLM activations, yielding improved stability and generalization over one-shot weight prediction.
    % \item (Optional) We augment activations with \textbf{behavioral probe fingerprints} and provide ablations showing when auxiliary probes improve cross-recipe generalization.
    \item We evaluate on diverse fine-tuning settings and unseen checkpoints, demonstrating recovery of safety behavior with minimal impact on downstream task performance.
\end{itemize}


% https://aclanthology.org/2023.findings-emnlp.808.pdf
\section{Related Works}
\label{sec:related_works}

\paragraph{Safety Degradation under Fine-Tuning.}
\citet{qi2024finetuning} show that fine-tuning on benign data compromises safety. \citet{qi2025tokensdeep} find alignment is shallow, encoded in only a few token positions. Defenses include LISA~\cite{huang2024lisa} (lazy alignment during training), Safe LoRA~\cite{hsu2024safe} (projection of LoRA updates), and Vaccine~\cite{yi2024vaccine} (perturbation-aware alignment). All require per-checkpoint intervention.

\paragraph{Safety Classifiers and Guardrails.}
Llama Guard~\cite{inan2023llama} and WildGuard~\cite{han2024wildguard} classify prompts independently of the target model.
HiddenDetect~\cite{jiang2025hiddendetect} monitors hidden states but requires per-model probe fitting.
These methods are either model-agnostic (ignoring fine-tuning effects) or per-model (not transferable).

\paragraph{Hypernetworks and Side Networks.}
Hypernetworks~\cite{ha2017hypernetworks, chauhan2024brief} generate target-network weights conditioned on context, with applications in continual learning~\cite{vonoswald2020continual} and multi-task adaptation~\cite{navon2023equivariant}. Recent work has scaled this paradigm to LLMs: \citet{liang2025dragdrop} generate LoRA parameters directly from task descriptions, achieving zero-shot adaptation without gradient-based fine-tuning. Our work pursues a complementary direction---rather than generating task-adaptation weights from text prompts, we generate \textit{safety} weights from model activations, targeting the orthogonal problem of safety restoration.
Ladder Side-Tuning (LST)~\cite{sung2022lst} attaches a small side network to a frozen backbone for parameter-efficient transfer. We are, to our knowledge, the first to combine hypernetworks with side-network architectures for safety restoration.
\section{Methodology}
\label{sec:method}

\begin{figure}[t]
  \includegraphics[width=\columnwidth]{example-image-golden}
  \caption{A figure with a caption that runs for more than one line.
    Example image is usually available through the \texttt{mwe} package
    without even mentioning it in the preamble.}
  \label{fig:experiments}
\end{figure}

\begin{figure*}[t]
  \includegraphics[width=0.48\linewidth]{example-image-a} \hfill
  \includegraphics[width=0.48\linewidth]{example-image-b}
  \caption {A minimal working example to demonstrate how to place
    two images side-by-side.}
\end{figure*}
\input{tables/example1}
\section{Experiments}
\label{sec:experiments}

% \subsection{Setup}

\mysection{Models and Domains.} We evaluate on Qwen2-7B-Instruct~\cite{yang2024qwen2} and LLaMA-3-8B-Instruct~\cite{dubey2024llama} across 22 datasets spanning QA, commonsense reasoning, math, knowledge, classification, code, and instruction-following (Appendix~\ref{app:domains}). Each dataset is fine-tuned with LoRA (rank 16, $\alpha{=}32$) and merged.

\paragraph{Safety Benchmarks.}
We evaluate on three harmful-prompt benchmarks: \textbf{BeaverTails}~\cite{ji2024beavertails} (700 prompts, 14 categories) is used as the harmful data source during \ssn{} training (\S\ref{sec:gt}); the evaluation split is disjoint from training. \textbf{AdvBench}~\cite{zou2023universal} (520 prompts) and \textbf{HEx-PHI}~\cite{qi2024finetuning} (330 prompts, 11 categories) are never seen in any training phase, serving as zero-shot safety tests.

\paragraph{Evaluation Protocol.}
We use three holdout splits (Table~\ref{tab:splits}), each removing four datasets from different task categories. The hypernetwork trains on the remaining 18 datasets and is evaluated on held-out datasets it has never seen. We report \textbf{Task Accuracy} (domain-specific test performance; official HuggingFace validation/test splits where available, 20\% holdout otherwise) and \textbf{Harmful Rate} (percentage of harmful prompts producing harmful output, lower is safer; 700 BeaverTails prompts disjoint from training, judged by the beaver-dam-7b moderation model~\cite{ji2024beavertails}). We compare against the \textbf{base model} (aligned, no fine-tuning) and the \textbf{LoRA model} (fine-tuned, no safety mechanism). Table~\ref{tab:comparison} provides a qualitative comparison with prior defense methods.

\begin{table}[h]
\caption{Holdout splits. Each split removes four datasets from different task categories. $^\dagger$20\% holdout eval (no standard test set).}
\label{tab:splits}
\centering
\renewcommand{\arraystretch}{1.4}
\resizebox{\columnwidth}{!}{%
\begin{tabular}{@{}cllll@{}}
\toprule
\textbf{Split} & \textbf{Domain 1} & \textbf{Domain 2} & \textbf{Domain 3} & \textbf{Domain 4} \\
\midrule
A & BoolQ {\small(QA)} & ARC {\small(Science)} & Comp.\ Math {\small(Math)} & Alpaca$^\dagger$ {\small(Instr.)} \\[3pt]
B & PIQA {\small(Physical)} & HellaSwag {\small(Common.)} & CSenseQA {\small(Common.)} & OpenHermes$^\dagger$ {\small(Instr.)} \\[3pt]
C & MMLU {\small(Knowledge)} & NQ-Open {\small(QA)} & AG News$^\dagger$ {\small(Classif.)} & Dolly$^\dagger$ {\small(Instr.)} \\
\bottomrule
\end{tabular}%
}
\end{table}

\subsection{Main Results}

Table~\ref{tab:main_results} shows results across all three splits. LoRA fine-tuning increases the harmful rate from 3--7\% (base) to as high as 43\%, with the largest degradation on instruction-following and commonsense domains. \method{} reduces the harmful rate to below 1\% across all 12 holdout datasets and both architectures, while preserving task accuracy within 1 percentage point. The highest residual rates (NQ-Open, 0.6\%) still represent $>$95\% relative reduction from the fine-tuned model.

% Combined main results — LLaMA-3-8B + Qwen2-7B, LoRA vs OneShot vs HyperSafe.
\begin{table}[t]
\centering
\renewcommand{\arraystretch}{1.1}
\resizebox{\textwidth}{!}{%
\begin{tabular}{@{}l ccc ccc ccc ccc@{}}
\toprule
& \multicolumn{3}{c}{\textbf{Task Acc (\%) $\uparrow$}}
& \multicolumn{3}{c}{\textbf{LoRA Harm (BT/AB/HX, \%) $\downarrow$}}
& \multicolumn{3}{c}{\textbf{OneShot Harm $\downarrow$}}
& \multicolumn{3}{c}{\textbf{\method{} Harm $\downarrow$}} \\
\cmidrule(lr){2-4} \cmidrule(lr){5-7} \cmidrule(lr){8-10} \cmidrule(lr){11-13}
\textbf{Domain} & LoRA & OneShot & \method{} & BT & AB & HX & BT & AB & HX & BT & AB & HX \\
\midrule
\multicolumn{13}{@{}l}{\textit{\textbf{LLaMA-3-8B}}} \\
\midrule
BoolQ                  & 87.2 & 87.2 & 86.9 &  9.7 &  1.5 &  7.0 & 0.3 & 0.0 & 1.0 & 0.1 & 0.0 & 0.0 \\
ARC                    & 86.6 & 18.8 & 86.5 & 41.7 & 33.8 & 20.0 & 0.1 & 0.0 & 0.7 & 0.0 & 0.0 & 0.0 \\
Comp.\ Math$^\dagger$  & 33.8 & 28.4 & 34.1 &  8.6 &  1.7 &  6.0 & 0.9 & 0.4 & 1.7 & 0.0 & 0.0 & 0.0 \\
Alpaca$^\dagger$       & 39.0 & 15.2 & 38.9 & 23.0 &  9.6 & 34.3 & 0.1 & 0.0 & 0.7 & 0.1 & 0.0 & 0.0 \\
PIQA                   & 88.1 & 87.2 & 87.7 & 37.7 & 40.4 & 37.7 & 0.4 & 0.0 & 0.7 & 0.0 & 0.0 & 0.0 \\
HellaSwag              & 89.4 & 87.0 & 89.1 & 12.1 & 15.2 & 11.3 & 0.4 & 0.4 & 0.7 & 0.1 & 0.0 & 0.3 \\
CSenseQA               & 81.7 & 79.9 & 81.5 & 17.7 &  5.4 &  5.0 & 0.0 & 1.9 & 6.7 & 0.0 & 0.0 & 0.0 \\
OpenHermes$^\dagger$   & 54.9 & 16.4 & 54.6 & 41.0 & 85.6 & 62.0 & 0.0 & 0.0 & 2.7 & 0.0 & 0.0 & 0.0 \\
MMLU                   & 61.1 & 61.5 & 61.0 &  9.9 & 10.2 &  8.7 & 0.3 & 0.0 & 1.0 & 0.0 & 0.0 & 0.0 \\
NQ-Open$^\dagger$      & 77.4 & 55.9 & 77.1 & 20.1 & 11.5 & 19.7 & 2.7 & 0.2 & 2.0 & 0.4 & 0.0 & 0.0 \\
AG News$^\dagger$      & 94.1 & 93.9 & 90.6 &  9.0 &  0.8 &  5.3 & 2.7 & 2.3 & 9.3 & 0.0 & 0.0 & 0.0 \\
Dolly$^\dagger$        & 37.8 & 16.9 & 34.4 & 38.6 & 64.0 & 49.0 & 0.1 & 0.0 & 0.7 & 0.1 & 0.0 & 0.0 \\
\textit{Avg.}          & \textit{69.3} & \textit{54.0} & \textit{68.5}
                       & \textit{22.4} & \textit{23.3} & \textit{22.2}
                       & \textit{0.7}  & \textit{0.4}  & \textit{2.3}
                       & \textit{0.1}  & \textit{0.0}  & \textit{0.0} \\
\midrule
\multicolumn{13}{@{}l}{\textit{\textbf{Qwen2-7B}}} \\
\midrule
BoolQ                  & 88.4 & 88.5 & 88.2 & 14.3 &  7.5 & 13.3 &  4.4 & 0.2 &  2.3 & 0.1 & 0.0 & 0.0 \\
ARC                    & 91.3 & 91.5 & 91.3 &  7.1 &  0.2 & 33.7 &  1.0 & 0.0 &  3.7 & 0.0 & 0.0 & 0.0 \\
Comp.\ Math$^\dagger$  & 74.4 & 76.1 & 74.1 &  6.0 &  0.0 &  6.7 &  5.4 & 0.0 &  8.3 & 0.0 & 0.0 & 0.0 \\
Alpaca$^\dagger$       & 39.6 & 37.2 & 39.5 & 19.4 & 13.3 & 33.0 &  5.4 & 1.0 &  8.0 & 0.4 & 0.0 & 0.0 \\
PIQA                   & 90.8 & 90.6 & 90.6 & 43.1 & 67.7 & 57.3 &  4.3 & 0.2 &  2.0 & 0.0 & 0.0 & 0.0 \\
HellaSwag              & 93.2 & 58.0 & 92.8 & 32.9 & 29.8 & 45.3 &  4.7 & 1.0 &  5.0 & 0.1 & 0.0 & 0.3 \\
CSenseQA               & 84.6 & 84.2 & 84.5 & 15.7 &  3.7 & 14.0 &  1.0 & 0.0 &  2.0 & 0.0 & 0.0 & 0.0 \\
OpenHermes$^\dagger$   & 57.2 & 54.6 & 57.0 & 38.0 & 70.2 & 64.7 &  1.1 & 3.3 &  8.3 & 0.1 & 0.0 & 0.1 \\
MMLU                   & 68.7 & 68.7 & 68.4 &  9.7 &  0.2 & 22.0 &  5.1 & 0.2 &  6.3 & 0.1 & 0.0 & 0.3 \\
NQ-Open$^\dagger$      & 54.8 & 23.7 & 54.3 & 11.9 &  6.0 & 28.7 & 11.9 & 1.0 & 14.7 & 0.6 & 0.0 & 0.7 \\
AG News$^\dagger$      & 97.6 & 92.5 & 97.6 & 10.1 &  0.2 &  6.0 &  3.7 & 0.0 &  6.7 & 0.0 & 0.0 & 0.0 \\
Dolly$^\dagger$        & 35.6 & 33.8 & 35.4 & 22.1 & 24.0 & 44.3 & 11.0 & 2.7 & 11.7 & 0.3 & 0.0 & 0.0 \\
\textit{Avg.}          & \textit{73.0} & \textit{66.6} & \textit{72.8}
                       & \textit{19.2} & \textit{18.6} & \textit{30.8}
                       & \textit{4.9}  & \textit{0.8}  & \textit{6.6}
                       & \textit{0.1}  & \textit{0.0}  & \textit{0.1} \\
\bottomrule
\end{tabular}%
}
\caption{Main results across 12 holdout domains on LLaMA-3-8B and Qwen2-7B. \method{} reduces harmful rates to near-zero on all three benchmarks (BeaverTails, AdvBench, HEx-PHI) while preserving accuracy within ${<}1$~percentage~point of the LoRA fine-tuned baseline. OneShot~\citep{kim2024oneshot} achieves comparable safety on most LLaMA-3 domains but catastrophically degrades accuracy on instruction-following datasets (ARC, Alpaca, OpenHermes, Dolly); on Qwen2 it leaves substantial residual harm (avg 6.6\% HEx-PHI). $^\dagger$20\% holdout eval (no standard test set).}
\label{tab:main_results}
\end{table}


\paragraph{Comparison with Llama Guard 4.}
We compare against Llama Guard 4 (12B), the latest publicly released version of the Llama Guard family (originally introduced by~\citet{inan2023llama}; Llama Guard 4 released as part of the Llama 4 suite~\citep{meta2025llama4}), on the same three benchmarks (Appendix Table~\ref{tab:llamaguard}). On the two fully zero-shot benchmarks --- the appropriate comparison, since BeaverTails appears in \ssn{} training (with a disjoint eval split) --- Llama Guard 4 misses 6.5\% of AdvBench and 6.7\% of HEx-PHI prompts, while \method{} misses 0.0\% and 0.3\% respectively. This is achieved with an \ssn{} of only $\sim$3\% of backbone parameters versus Llama Guard 4's dedicated 12B model. The gap is wider on BeaverTails (53.6\% vs.\ 0.7\%).

\paragraph{Ablations and additional analyses.}
Five ablations are reported in full in Appendices~\ref{app:ablation_summary}--\ref{app:additional}: (i)~loss components, (ii)~training-data diversity (instruction-following is the most important category), (iii)~calibration set size (saturates at $N_\text{cal}{=}50$), (iv)~generalization to full fine-tuning (\ssn{}s trained for LoRA models transfer), and (v)~cross-benchmark safety transfer ($<$1\% harm on the zero-shot AdvBench and HEx-PHI benchmarks).
\section{Conclusion and Limitations}
\label{sec:conclusion}

\method{} generates model-specific \ssn{}s that reduce harmful outputs to near zero on held-out checkpoints of two architectures while preserving task accuracy, and fingerprint-conditioned generation contributes substantially to these gains over a shared \ssn{}. For a provider maintaining many checkpoints of one backbone, the one-time meta-training is amortized across all of them.

\mysection{Limitations.}
\method{} is designed as an input-side defense against safety degradation from benign fine-tuning, rather than as a general jailbreak defense. White-box adaptive attacks on the generated \ssn{}, harms that emerge across turns, and output-side failures are outside our evaluation; checkpoint-specific weights alone do not provide security, and extending the \ssn{} to score generated text is a natural next step. \method{} assumes white-box access to activations and a backbone matching the hypernetwork, whose one-time meta-training uses safety data and is repeated per backbone family. Finally, our evaluation relies on two automated judges without human validation, measures over-refusal on XSTest and OR-Bench, evaluates baselines at a single operating point, reports seed variance for hypernetwork training, and covers 7--8B models; broadening each of these is left to future work.


%\begin{ack}
%Use unnumbered first level headings for the acknowledgments. All acknowledgments
%go at the end of the paper before the list of references. Moreover, you are required to declare
%funding (financial activities supporting the submitted work) and competing interests (related financial activities outside the submitted work).
%More information about this disclosure can be found at: \url{https://neurips.cc/Conferences/2026/PaperInformation/FundingDisclosure}.
%
%
%Do {\bf not} include this section in the anonymized submission, only in the final paper. You can use the \texttt{ack} environment provided in the style file to automatically hide this section in the anonymized submission.
%\end{ack}

\bibliographystyle{unsrt}
\bibliography{custom}

%\section*{References}
%
%
%References follow the acknowledgments in the camera-ready paper. Use unnumbered first-level heading for
%the references. Any choice of citation style is acceptable as long as you are
%consistent. It is permissible to reduce the font size to \verb+small+ (9 point)
%when listing the references.
%Note that the Reference section does not count towards the page limit.
%\medskip
%
%
%{
%\small
%
%
%[1] Alexander, J.A.\ \& Mozer, M.C.\ (1995) Template-based algorithms for
%connectionist rule extraction. In G.\ Tesauro, D.S.\ Touretzky and T.K.\ Leen
%(eds.), {\it Advances in Neural Information Processing Systems 7},
%pp.\ 609--616. Cambridge, MA: MIT Press.
%
%
%[2] Bower, J.M.\ \& Beeman, D.\ (1995) {\it The Book of GENESIS: Exploring
%  Realistic Neural Models with the GEneral NEural SImulation System.}  New York:
%TELOS/Springer--Verlag.
%
%
%[3] Hasselmo, M.E., Schnell, E.\ \& Barkai, E.\ (1995) Dynamics of learning and
%recall at excitatory recurrent synapses and cholinergic modulation in rat
%hippocampal region CA3. {\it Journal of Neuroscience} {\bf 15}(7):5249-5262.
%}


%%%%%%%%%%%%%%%%%%%%%%%%%%%%%%%%%%%%%%%%%%%%%%%%%%%%%%%%%%%%

\clearpage
%\section{Example Appendix}
%\label{sec:appendix}
%
%This is an appendix.

\appendix

\section{Domain Dataset Details}
\label{app:domains}

\begin{table}[h]
\centering
\small
\begin{tabular}{@{}llr@{}}
\toprule
\textbf{Domain} & \textbf{Category} & \textbf{Train Size} \\
\midrule
BoolQ & QA & 9.4k \\
NQ-Open & QA & 10k \\
SQuAD & QA & 26k \\
PubMedQA & QA & 0.5k \\
ARC & Science QA & 2.4k \\
CommonsenseQA & Commonsense & 9.7k \\
PIQA & Physical reasoning & 16.1k \\
HellaSwag & Commonsense & 10k \\
SocialIQA & Social reasoning & 10k \\
GSM8K & Math & 7.5k \\
Comp.\ Math & Math & 7.5k \\
CodeFeedback & Code & 10k \\
AG News & Classification & 24k \\
MMLU & Knowledge & 14k \\
Platypus & Knowledge & 5k \\
Dolly & Instruction & 15k \\
OpenOrca & Instruction & 10k \\
Alpaca & Instruction & 52k \\
FLAN & Instruction & 10k \\
OpenHermes & Instruction & 10k \\
NoRobots & Instruction & 10k \\
WizardLM & Instruction & 10k \\
\bottomrule
\end{tabular}
\caption{The 22 fine-tuning domains. All fine-tuned with LoRA (rank 16, $\alpha{=}32$) and merged.}
\label{tab:domains}
\end{table}

\section{Implementation Details}
\label{app:implementation}

Table~\ref{tab:hyperparams} summarizes all hyperparameters.

\begin{table}[h]
\centering
\renewcommand{\arraystretch}{1.2}
\begin{tabular}{@{}llc@{}}
\toprule
\textbf{Component} & \textbf{Parameter} & \textbf{Value} \\
\midrule
\multirow{4}{*}{Side Network}
& Hidden dim ($h_s$)        & $h/4$ \\
& Layers ($K$)              & 12 \\
& Attention heads            & 7 / 8 \\
& Overhead                   & $\sim$3\% \\
\midrule
\multirow{4}{*}{LoRA Fine-Tuning}
& Rank                       & 16 \\
& $\alpha$                   & 32 \\
& Target modules             & all linear \\
& Epochs / lr                & 3 / $2{\times}10^{-4}$ \\
\midrule
\multirow{8}{*}{Hypernetwork}
& Hidden dim                 & 512 \\
& Classifier rank            & 16 \\
& Direction proj.\ dim       & 256 \\
& Magnitude hidden dim       & 128 \\
& $N_\text{cal}$             & 50 \\
& Epochs / lr                & 100 / $5{\times}10^{-5}$ \\
& $\alpha$ ($\mathcal{L}_\text{cls}$) & 10 \\
& $\beta$ ($\mathcal{L}_\text{func}$) & 5 \\
\bottomrule
\end{tabular}
\caption{Hyperparameters. Side network heads: 7 for Qwen2-7B ($h_s{=}896$), 8 for LLaMA-3-8B ($h_s{=}1024$).}
\label{tab:hyperparams}
\end{table}

\section{Additional Ablations}
\label{app:ablations}

\paragraph{Hypernetwork Architecture.}
Table~\ref{tab:ablation_arch} varies the generation rank and hidden dimension, reporting per-domain task accuracy on Split~A holdout. Lower capacity (rank~2, dim~256) degrades accuracy across all domains; increasing beyond the default yields marginal gains. Safety remains robust across all configurations.

\begin{table}[h]
\centering
\small
\resizebox{\columnwidth}{!}{%
\begin{tabular}{@{}lccccc@{}}
\toprule
& \multicolumn{4}{c}{\textbf{Task Acc (\%) $\uparrow$}} & \textbf{Harm.} \\
\cmidrule(lr){2-5}
\textbf{Config} & BoolQ & ARC & Math$^\dagger$ & Alpaca$^\dagger$ & \textbf{Rate $\downarrow$} \\
\midrule
Rank 2          & 84.3 & 85.8 & 33.5 & 37.1 & 0.2\% \\
Rank 4 (def.)   & 86.9 & 86.5 & 34.1 & 38.9 & 0.1\% \\
Rank 8          & 87.0 & 86.5 & 34.2 & 39.0 & 0.1\% \\
\addlinespace
Dim 256         & 85.4 & 85.1 & 33.3 & 37.4 & 0.3\% \\
Dim 512 (def.)  & 86.9 & 86.5 & 34.1 & 38.9 & 0.1\% \\
Dim 1024        & 87.1 & 86.6 & 34.1 & 39.0 & 0.1\% \\
\bottomrule
\end{tabular}%
}
\caption{Architecture sensitivity (LLaMA-3, Split~A holdout). Lower capacity degrades accuracy; higher capacity gives marginal gains. Safety is robust across all configurations.}
\label{tab:ablation_arch}
\end{table}

\section{Per-Domain Safety Scores}
\label{app:per_domain}

\begin{table}[H]
\centering
\small
\resizebox{\columnwidth}{!}{%
\begin{tabular}{@{}llccc@{}}
\toprule
& & \multicolumn{3}{c}{\textbf{Harmful Rate (\%) $\downarrow$}} \\
\cmidrule(lr){3-5}
\textbf{Model} & \textbf{Domain} & BT & AdvBench & HEx-PHI \\
\midrule
\multirow{6}{*}{\rotatebox{90}{\small LLaMA-3}}
& BoolQ       & 0.1 & 0.0 & 0.0 \\
& ARC         & 0.0 & 0.0 & 0.0 \\
& PIQA        & 0.0 & 0.0 & 0.0 \\
& HellaSwag   & 0.1 & 0.0 & 0.3 \\
& MMLU        & 0.0 & 0.0 & 0.0 \\
& NQ-Open     & 0.4 & 0.0 & 0.0 \\
\midrule
\multirow{6}{*}{\rotatebox{90}{\small Qwen2}}
& BoolQ       & 0.1 & 0.0 & 0.0 \\
& ARC         & 0.0 & 0.0 & 0.0 \\
& PIQA        & 0.0 & 0.0 & 0.0 \\
& HellaSwag   & 0.1 & 0.0 & 0.3 \\
& MMLU        & 0.1 & 0.0 & 0.3 \\
& NQ-Open     & 0.6 & 0.0 & 0.7 \\
\bottomrule
\end{tabular}%
}
\caption{Per-domain harmful rates across three benchmarks (domains with official HF eval splits). AdvBench and HEx-PHI achieve near-zero harmful rates across the board. BT eval split is disjoint from training.}
\label{tab:per_domain_safety}
\end{table}

%%%%%%%%%%%%%%%%%%%%%%%%%%%%%%%%%%%%%%%%%%%%%%%%%%%%%%%%%%%%

\newpage
\input{checklist.tex}


\end{document}